\subsection{Linear mappings between Lebesgue spaces defined by a kernel}

Here we take K = C. Let X be a locally compact topological space equipped with a positive measure \(\mu\) . For every \(r \in [1, +\infty]\) , we denote \(\mathcal{L}^r (\mathrm{X}) = \mathcal{L}_{\mathbf{C}}^r (\mathrm{X},\mu)\) and \(\mathrm{L}^r (\mathrm{X}) = \mathrm{L}_{\mathbf{C}}^r (\mathrm{X},\mu)\) . When \(r\in [1, + \infty [\) , one identifies the dual of \(\mathrm{L}^r (\mathrm{X})\) with \(\mathrm{L}^{r^{\prime}}(\mathrm{X})\) where \(r^\prime\) is the conjugate exponent of \(r\) such that \(1 / r + 1 / r' = 1\) (INT, V, p. 61, § 5, no. 8, th. 4). We denote by \(\| f\| _r\) the norm (or semi-norm) of \(f\in \mathrm{L}^r (\mathrm{X})\) (or \(f\in \mathcal{L}^r (\mathrm{X})\) ).

Recall that a subset A of X is said to be \(\mu\) -moderate (INT, V, p. 4, § 1, n° 2) if it is contained in the union of a sequence of \(\mu\) -integrable sets, and that a function \(f\) defined on X with values in a vector space or in \(\overline{\mathbf{R}}\) is \(\mu\) -moderate if it vanishes on the complement of a \(\mu\) -moderate subset of X. If \(1 \leqslant p < +\infty\) and \(f \in \mathcal{L}^p(\mathrm{X})\) , then \(f\) is moderate (INT, V, p. 9, § 1, no. 3, cor.). If \(f\) is \(\mu\) -moderate, then \(fg\) is \(\mu\) -moderate for every function \(g\) .

Let \(p\) and \(q\) be elements of \(]1, +\infty[\) . Let X and Y be locally compact topological spaces, equipped with positive measures \(\mu\) and \(\nu\) , respectively. One equips the space \(X \times Y\) (resp. \(Y \times X\) ) with the product measure \(\mu \otimes \nu\) (resp. \(\nu \otimes \mu\) ).

We denote \(\mathcal{N}^{p,q}(\mathrm{X}\times\mathrm{Y},\mu\otimes\nu)\) , or simply \(\mathcal{N}^{p,q}(\mathrm{X}\times\mathrm{Y})\) , the set of functions k measurable with respect to \((\mu\otimes\nu)\) from X × Y into C such that there exists a real number \(C\geqslant0\) satisfying

\[
\int_ {\mathrm{X} \times \mathrm{Y}} ^ {*} | k (x, y) f (x) g (y) | d (\mu \otimes \nu) (x, y) \leqslant \mathrm{C} \| f \| _ {p} \| g \| _ {q}\tag{2}
\]

for all functions \(f \in \mathcal{L}^p(\mathrm{X})\) and \(g \in \mathcal{L}^q(\mathrm{Y})\) . We then denote by \(\| k \|_{p,q}\) the lower bound of the set of real numbers C having this property.

Remark. --- Let \(f \in \mathcal{L}^{p}(\mathrm{X})\) and \(g \in \mathcal{L}^{q}(\mathrm{Y})\) . Since \(p < +\infty\) and \(q < +\infty\) , the functions f and g are moderate, and the function defined by \((x, y) \mapsto f(x)g(y)\) is \((\mu \otimes \nu)\) -moderate (cf. INT, V, p. 92, § 8, n°3, cor. 2). Consequently, the function \((x, y) \mapsto k(x, y)f(x)g(y)\) is \((\mu \otimes \nu)\) -moderate, and in particular the upper integral of its absolute value coincides with the essential upper integral (INT, V, p. 6, § 1, n°3, prop. 7).

The set \(\mathcal{N}^{p,q}(\mathrm{X}\times \mathrm{Y})\) is a vector subspace of the space of functions from \(\mathrm{X}\times\mathrm{Y}\) into \(\mathbf{C}\) and the mapping defined by \(k\mapsto \| k\|_{p,q}\) is a semi-norm on \(\mathcal{N}^{p,q}(\mathrm{X}\times \mathrm{Y})\) .

Lemma 1. --- Let \(k \in \mathcal{N}^{p,q}(X \times Y)\) . The function \(k\) is locally \((\mu \otimes \nu)\) -integrable.

The map \(k\) is \((\mu \otimes \nu)\) -measurable by hypothesis. Let A be a compact subset of \(\mathrm{X} \times \mathrm{Y}\) . There exist compact subsets \(\mathrm{B} \subset \mathrm{X}\) and C ⊂ Y such that A ⊂ B × C. Applying (2) to the characteristic functions \(\varphi_{B}\) of B and \(\varphi_{C}\) of C, we obtain

\[
\begin{array}{l} \int_ {X \times Y} ^ {*} | k (x, y) | \varphi_ {\mathrm{A}} (x, y) d (\mu \otimes \nu) (x, y) \\ \qquad \leqslant \int_ {X \times Y} ^ {*} | k (x, y) | \varphi_ {\mathrm{B}} (x) \varphi_ {\mathrm{C}} (y) d (\mu \otimes \nu) (x, y) <   + \infty , \end{array}
\]

whence the result (INT, V, p. 41, §5, n° 1, prop. 1 and def. 1).

PROPOSITION 3. --- a) Let \(k \in \mathcal{N}^{p,q}(X \times Y)\) . For \(f \in \mathcal{L}^p(X)\) and \(g \in \mathcal{L}^q(Y)\) , the function

\[
(x, y) \mapsto k (x, y) f (x) g (y)
\]

is \((\mu \otimes \nu)\) -integrable and there exists a unique map \(u_{k}\) from \(\mathrm{L}^{p}(\mathrm{X})\) into \(\mathrm{L}^{q'}(\mathrm{Y})\) such that

\[
\langle g, u _ {k} (f) \rangle = \int_ {\mathrm{X} \times \mathrm{Y}} k (x, y) f (x) g (y) d (\mu \otimes \nu) (x, y)
\]

for all \(f \in \mathrm{L}^p(\mathrm{X})\) and \(g \in \mathrm{L}^q(\mathrm{Y})\) . The map \(u_k\) is linear and continuous, its norm is \(\leqslant \| k\|_{p,q}\) ;

\begin{enumerate}
\def\labelenumi{\alph{enumi})}
\setcounter{enumi}{1}
\item
  Let \(k \in \mathcal{N}^{p,q}(X \times Y)\) . We have \(u_k = 0\) if and only if k is locally \((\mu \otimes \nu)\) -negligible;
\item
  The map which to k associates \(u_{k}\) is a continuous linear map from the semi-normed space \(\mathcal{N}^{p,q}(\mathrm{X}\times\mathrm{Y})\) to the Banach space \(\mathcal{L}(\mathrm{L}^{p}(\mathrm{X});\mathrm{L}^{q^{\prime}}(\mathrm{Y}))\) , such that \(\|u_{k}\|\leqslant\|k\|_{p,q}\) for all \(k\in\mathcal{N}^{p,q}(\mathrm{X}\times\mathrm{Y})\) .
\end{enumerate}

Let \(k \in \mathcal{N}^{p,q}(\mathrm{X} \times \mathrm{X})\) . By definition, the function from \(\mathrm{X} \times \mathrm{Y}\) into \(\mathbf{C}\) defined by \((x,y) \mapsto k(x,y)f(x)g(y)\) is \((\mu \otimes \nu)\) -integrable when \(f \in \mathcal{L}^p(\mathrm{X})\) and \(g \in \mathcal{L}^q(\mathrm{Y})\) , and the map

\[
b _ {k} \colon (f, g) \mapsto \int_ {\mathrm{X} \times \mathrm{Y}} k (x, y) f (x) g (y) d (\mu \otimes \nu) (x, y)
\]

is a continuous bilinear map from \(\mathrm{L}^{p}(\mathrm{X})\times\mathrm{L}^{q}(\mathrm{Y})\) to C. For every \(f\in\mathrm{L}^{p}(\mathrm{X})\) , there therefore exists a unique \(h\in\mathrm{L}^{q^{\prime}}(\mathrm{Y})\) such that \(b_{k}(f,g)=\langle g,h\rangle\) for all \(g\in\mathrm{L}^{q}(\mathrm{Y})\) ; denote it \(u_{k}(f)=h\) . By (2), we have \(\|u_{k}(f)\|_{q^{\prime}}\leqslant\|k\|_{p,q}\|f\|_{p}\) ; the map \(u_{k}\) is linear and continuous with norm \(\leqslant\|k\|_{p,q}\) . Moreover, the map \(k\mapsto u_{k}\) is linear, and assertions a) and c) follow from this.

Let us prove \(b\) ). Suppose that \(k\) is locally \((\mu \otimes \nu)\) -negligible. Let \(f \in \mathcal{L}^p(\mathrm{X})\) and \(g \in \mathcal{L}^q(\mathrm{Y})\) ; the function defined on \(\mathrm{X} \times \mathrm{Y}\) by \((x, y) \mapsto k(x, y)f(x)g(y)\) is \((\mu \otimes \nu)\) -moderate. Since it is locally \((\mu \otimes \nu)\) -negligible, it is negligible (INT, V, p. 7, § 1, n°2, cor. 1). Thus, we have \(\langle u_k(f), g \rangle = b_k(f, g) = 0\) . It follows that \(u_k = 0\) .

Conversely, let \(k \in \mathcal{N}^{p,q}(X \times Y)\) be such that \(u_k = 0\) . The function k is locally \((\mu \otimes \nu)\) -integrable (Lemma 1). For all functions \(f \in \mathcal{K}(X)\) and \(g \in \mathcal{K}(Y)\) , we have

\[
\langle g, u _ {k} (f) \rangle = \int_ {\mathrm{X} \times \mathrm{Y}} f (x) g (y) k (x, y) d (\mu \otimes \nu) (x, y),
\]

therefore the measure \(k \cdot (\mu \otimes \nu)\) on X × Y is zero (cf. INT, III, p. 82, § 4, no. 1, Th. 1). This means that k is locally \((\mu \otimes \nu)\) -negligible (INT, V, p. 46, §5, no. 3, Cor. 2).

DEFINITION 1. --- Let \(k \in \mathcal{N}^{p,q}(\mathrm{X} \times \mathrm{Y})\) . We say that the map \(u_k\) is the integral operator with kernel \(k\) from \(\mathrm{L}^p(\mathrm{X})\) into \(\mathrm{L}^{q'}(\mathrm{Y})\) .

For \(k \in \mathcal{N}^{p,q}(\mathrm{X} \times \mathrm{Y})\) and \(f \in \mathcal{L}^p(\mathrm{X})\) , we will sometimes also denote by \(u_k(f)\) the image under \(u_k\) of the class of \(f\) in \(\mathrm{L}^p(\mathrm{X})\) .

PROPOSITION 4. --- Let \(k \in \mathcal{N}^{p,q}(\mathrm{X} \times \mathrm{Y})\) and \(f \in \mathcal{L}^p(\mathrm{X})\) . The mapping defined on \(\mathrm{X}\) by \(x \mapsto k(x,y)f(x)\) is \(\mu\) -integrable for every \(y\) outside a locally \(\nu\) -negligible set, and \(u_k(f)\) coincides with the class of the function

\[
y \mapsto \int_ {\mathrm{X}} k (x, y) f (x) d \mu (x)
\]

defined locally ν-almost everywhere.

Let \(Y'\) denote the set of \(y \in Y\) such that the mapping \(x \mapsto k(x, y)f(x)\) is not \(\mu\) -integrable. Let \(y \in Y\) and let C be a compact neighborhood of y, whose characteristic function we denote by \(\varphi\) . By condition (2) applied to f and \(\varphi\) , the function \((x, y) \mapsto k(x, y)f(x)\varphi(y)\) is integrable on X × Y with respect to the measure \(\mu \otimes \nu\) . By the Lebesgue--Fubini theorem (INT, V, p. 96, § 8, n°4, th. 1, a)), the function \(x \mapsto k(x, y)f(x)\varphi(y)\) is \(\mu\) -integrable for y outside a \(\nu\) -negligible set \(Y_{C}\) . Since \(Y' \cap C \subset Y_{C}\) , it follows that \(Y'\) is locally \(\nu\) -negligible in Y (INT, IV, p. 172, § 5, n°2, def. 3).

Let h denote the mapping defined locally \(\nu\) -almost everywhere on Y by \(y \mapsto \int k(x, y)f(x)d\mu\) . It is locally \(\nu\) -integrable (INT, V, loc.

cit.). Let \(g \in \mathcal{K}(\mathrm{Y})\) . We have

\[
\begin{array}{l} \langle g, u _ {k} (f) \rangle = \int_ {\mathrm{X} \times \mathrm{Y}} k (x, y) f (x) g (y) d (\mu \otimes \nu) (x, y) \\ \qquad = \int_ {\mathrm{Y}} \Bigl (\int_ {\mathrm{X}} k (x, y) f (x) g (y) d \mu (x) \Bigr) d \nu (y) \\ \qquad = \int_ {\mathrm{Y}} g (y) h (y) d \nu (y) \end{array}
\]

(INT, V, loc. cit.). We therefore have \(u_{k}(f) = h\) locally \(\nu\) -almost everywhere (cf. INT, V, p. 46, § 5, n°3, cor. 2).

PROPOSITION 5. --- The space \(\mathcal{L}^{p'}(\mathrm{X}\times\mathrm{Y})\) is contained in \(\mathcal{N}^{p,p}(\mathrm{X}\times\mathrm{Y})\) . The linear map from \(\mathcal{L}^{p'}(\mathrm{X}\times\mathrm{Y})\) into \(\mathcal{L}(\mathrm{L}^{p}(\mathrm{X});\mathrm{L}^{p'}(\mathrm{Y}))\) defined by \(k\mapsto u_{k}\) induces, by passing to the quotient, a continuous injective linear map from \(\mathrm{L}^{p'}(\mathrm{X}\times\mathrm{Y})\) into \(\mathcal{L}(\mathrm{L}^{p}(\mathrm{X});\mathrm{L}^{p'}(\mathrm{Y}))\) .

Let \(k \in \mathcal{L}^{p'}(\mathrm{X} \times \mathrm{Y})\) . The function \(k\) is measurable (INT, IV, p. 84, § 5, n° 6, th. 5). Let \(f \in \mathcal{L}^p(\mathrm{X})\) and \(g \in \mathcal{L}^p(\mathrm{Y})\) . The function \(h \colon (x,y) \mapsto f(x)g(y)\) belongs to \(\mathcal{L}^p(\mathrm{X} \times \mathrm{Y})\) and \(\|h\|_p = \|f\|_p \|g\|_p\) (INT, V, p. 95, § 8, n° 3, cor. 2). The function \(hk\) is therefore integrable on \(\mathrm{X} \times \mathrm{Y}\) and we have

\[
(\mu \otimes \nu) ^ {*} (| h k |) \leqslant \| k \| _ {p ^ {\prime}} \| h \| _ {p} = \| k \| _ {p ^ {\prime}} \| f \| _ {p} \| g \| _ {p}
\]

by Hölder's inequality (INT, IV, p. 208, § 6, n° 4, th. 2), which proves that \(k \in \mathcal{N}^{p,p}(\mathrm{X} \times \mathrm{Y})\) with \(\| k \|_{p,p} \leqslant \| k \|_{p'}\) .

Since \(p' \neq +\infty\) , every function \(k \in \mathcal{L}^{p'}(\mathrm{X} \times \mathrm{Y})\) is moderate (INT, V, p. 9, § 1, n° 3, cor.), and the last assertion then follows from prop. 3, b) and c).

For any function \(k\colon X \times Y \to \mathbf{C}\) and for every \(y \in Y\) , we denote by \(k_y\) the mapping from X into \(\mathbf{C}\) given by \(k_y(x) = k(x, y)\) , and we denote by \(k_\circ\) the mapping from Y into \(\mathcal{F}(X; \mathbf{C})\) defined by \(y \mapsto k_y\) . We shall likewise denote the mapping from Y into the space of functions defined \(\mu\) -almost everywhere on X that associates with \(y\) the class of \(k_y\) .

For any function \(k\) from \(\mathrm{X} \times \mathrm{Y}\) into \(\mathbf{C}\) , we denote

\[
\mathrm{N} _ {p ^ {\prime}, q ^ {\prime}} (k) = \left(\int_ {\mathrm{Y}} ^ {\bullet} \mathrm{N} _ {p ^ {\prime}} (k _ {y}) ^ {q ^ {\prime}} d \nu (y)\right) ^ {1 / q ^ {\prime}}.
\]

This is an element of \([0, +\infty]\) .

Lemma 2. --- For every \(k \in \mathcal{L}^{p'}(\mathrm{X} \times \mathrm{Y})\) , we have \(\mathrm{N}_{p', p'}(k) = \| k \|_{p'}\) .

Let \(k \in \mathcal{L}^{p'}(\mathrm{X} \times \mathrm{Y})\) . By INT, V, p. 96, § 8, no. 4, th. 1, a), the set \(Y'\) of elements \(y \in \mathrm{Y}\) such that the function \(|k_y|^{p'}\) is not \(\mu\) -integrable is \(\nu\) -negligible; moreover, the function defined by \(y \mapsto \mu(|k_y|^{p'})\) on \(\mathrm{Y} - \mathrm{Y}'\) and equal to zero on \(\mathrm{Y}'\) is \(\nu\) -integrable and satisfies

\[
\| k \| _ {p ^ {\prime}} ^ {p ^ {\prime}} = \int_ {\mathrm{X} \times \mathrm{Y}} | k (x, y) | ^ {p ^ {\prime}} d (\mu \otimes \nu) (x, y) = \int_ {\mathrm{Y}} \mu (| k _ {y} | ^ {p ^ {\prime}}) d \nu = \mathrm{N} _ {p ^ {\prime}, p ^ {\prime}} (k) ^ {p ^ {\prime}}.
\]

We denote \(\mathcal{L}^{p',q'}(\mathrm{X},\mathrm{Y},\mu,\nu)\) , or simply \(\mathcal{L}^{p',q'}(\mathrm{X},\mathrm{Y})\) , the complex vector space of functions \((\mu\otimes\nu)\) -measurable \(k\) from \(\mathrm{X}\times\mathrm{Y}\) into \(\mathbf{C}\) such that, for \(\nu\) -almost all \(y\in\mathrm{Y}\) , the map \(k_y\) belongs to \(\mathcal{L}^{p'}(\mathrm{X})\) , and such that the map \(k_\circ\) belongs to \(\mathcal{L}^{q'}_{\mathrm{L}^{p'}(\mathrm{X})}(\mathrm{Y},\nu)\) (INT, IV, p. 129, § 3, no. 4, def. 2). We equip \(\mathcal{L}^{p',q'}(\mathrm{X},\mathrm{Y})\) with the seminorm \(k\mapsto\mathrm{N}_{p',q'}(k)\) ; note that this then satisfies \(\mathrm{N}_{p',q'}(k)=\|k_\circ\|_{q'}\) , where \(k_\circ\) is viewed as a mapping with values in the Banach space \(\mathrm{L}^{p'}(\mathrm{X})\) .

PROPOSITION 6. --- a) The space \(\mathcal{L}^{p'}(\mathrm{X} \times \mathrm{Y})\) is contained in \(\mathcal{L}^{p',p'}(\mathrm{X},\mathrm{Y})\) and for \(k \in \mathcal{L}^{p'}(\mathrm{X} \times \mathrm{Y})\) , we have \(\mathrm{N}_{p',p'}(k) = \| k\|_{p'}\) ;

\begin{enumerate}
\def\labelenumi{\alph{enumi})}
\setcounter{enumi}{1}
\tightlist
\item
  The space \(\mathcal{L}^{p',q'}(\mathrm{X},\mathrm{Y})\) is contained in \(\mathcal{N}^{p,q}(\mathrm{X}\times \mathrm{Y})\) and for \(k\in \mathcal{L}^{p',q'}(\mathrm{X},\mathrm{Y})\) , we have \(\| k\|_{p,q}\leqslant \mathrm{N}_{p',q'}(k)\) .
\end{enumerate}

Let us prove assertion \(a\) ). Let \(k \in \mathcal{L}^{p'}(\mathrm{X} \times \mathrm{Y})\) . By INT, V, p. 96, § 8, no. 4, th. 1, \(a\) ), the function \(k_y\) belongs to \(\mathcal{L}^{p'}(\mathrm{X})\) for \(\nu\) -almost all \(y \in \mathrm{Y}\) .

Let \(\varepsilon > 0\) . There exists \(h \in \mathcal{K}(\mathrm{X} \times \mathrm{Y})\) such that \(\|k - h\|_{p'} < \varepsilon\) . The function \(h_{\circ}\) is a continuous function with compact support from Y into \(\mathrm{L}^{p'}(\mathrm{X})\) satisfying

\[
\| k _ {\circ} - h _ {\circ} \| _ {p ^ {\prime}} = \mathrm{N} _ {p ^ {\prime}, p ^ {\prime}} (k - h) = \| k - h \| _ {p ^ {\prime}} <   \varepsilon
\]

by Lemma 2. It follows that \(k_{\circ}\) belongs to \(\mathcal{L}_{\mathrm{L}^{p^{\prime}}(\mathrm{X})}^{p^{\prime}}(\mathrm{Y},\nu)\) , that is, \(k\) belongs to \(\mathcal{L}^{p^{\prime},p^{\prime}}(\mathrm{X},\mathrm{Y})\) . Lemma 2 then shows that \(\mathrm{N}_{p^{\prime},p^{\prime}}(k)=\|k\|_{p^{\prime}}\) .

Let us prove assertion \(b\) ). Let \(k \in \mathcal{L}^{p', q'}(\mathrm{X}, \mathrm{Y})\) and let \(f \in \mathcal{L}^p(\mathrm{X})\) . For every \(y \in \mathrm{Y}\) , Hölder's inequality (INT, IV, p. 208, § 6, no. 4, th. 2) applied to \(k_y\) implies

\[
\int_ {X} ^ {*} | k (x, y) f (x) | d \mu (x) \leqslant \| k _ {y} \| _ {p ^ {\prime}} \| f \| _ {p}.
\]

Let \(g \in \mathcal{L}^{q'}(\mathrm{Y})\) . The maps \(f\) and \(g\) are moderate, since \(p\) and \(q'\) are finite. Consequently, the map \((x,y) \mapsto f(x)g(y)\) is \((\mu \otimes \nu)\) -moderate (cf. INT, V, p. 92, § 8, no. 3, cor. 2), and hence the map \((x,y)\mapsto k(x,y)f(x)g(y)\) is \((\mu\otimes\nu)\) -moderate. According to INT, V, p. 93, § 8, no. 3, prop. 7, a), it then follows that

\[
\begin{array}{l} \int_ {\mathrm{X} \times \mathrm{Y}} ^ {*} | k (x, y) f (x) g (y) | d (\mu \otimes \nu) \\ = \int_ {\mathrm{Y}} ^ {*} | g (y) | \Big (\int_ {\mathrm{X}} ^ {*} | k (x, y) f (x) | d \mu (x) \Big) d \nu (y) \\ \leqslant \| f \| _ {p} \int_ {\mathrm{Y}} ^ {*} \| k _ {y} \| _ {p ^ {\prime}} | g (y) | d \nu (y) \leqslant \mathrm{N} _ {p ^ {\prime}, q ^ {\prime}} (k) \| f \| _ {p} \| g \| _ {q}, \end{array}
\]

using Hölder's inequality once again. This concludes the proof.

PROPOSITION 7. --- Let \(k \in \mathcal{L}^{p',q'}(\mathrm{X},\mathrm{Y})\) . The linear map \(u_k\) from \(\mathrm{L}^p (\mathrm{X})\) into \(\mathrm{L}^{q'}(\mathrm{Y})\) is compact.

Suppose first that \(k\) is such that \(k_{\circ} \in \mathcal{K}(\mathrm{Y}; \mathcal{L}^{p'}(\mathrm{X}))\) and denote by A the support of \(k_{\circ}\) . Let F be the subspace of \(\mathcal{K}(\mathrm{Y}, \mathrm{A}; \mathcal{L}^{p'}(\mathrm{X}, \mu))\) generated by the functions \(y \mapsto f_2(y)f_1\) , where \(f_2 \in \mathcal{K}(\mathrm{Y}, \mathrm{A})\) and \(f_1 \in \mathcal{L}^{p'}(\mathrm{X})\) . When \(m \in \mathcal{K}(\mathrm{Y}; \mathcal{L}^{p'}(\mathrm{X}))\) satisfies \(m_{\circ} \in \mathrm{F}\) , the linear map \(u_m\) is of finite rank, hence compact. For \(m_1\) and \(m_2\) in \(\mathcal{K}(\mathrm{Y}, \mathrm{A}; \mathcal{L}^{p'}(\mathrm{X}))\) we also have

\[
\begin{array}{c} \| u _ {m _ {1}} - u _ {m _ {2}} \| \leqslant \| m _ {1} - m _ {2} \| _ {p ^ {\prime}, q ^ {\prime}} = \Big (\int_ {\mathrm{Y}} \| m _ {1, y} - m _ {2, y} \| _ {p ^ {\prime}} ^ {q ^ {\prime}} d \nu (y) \Big) ^ {1 / q ^ {\prime}} \\ \leqslant \nu (\mathrm{A}) ^ {1 / q ^ {\prime}} \sup _ {y \in \mathrm{Y}} \mathrm{N} _ {p ^ {\prime}} (m _ {1, y} - m _ {2, y}). \end{array}
\]

Since the space F is dense in \(\mathcal{K}(\mathrm{Y},\mathrm{A};\mathrm{L}^{p^{\prime}}(\mathrm{X}))\) for the topology of uniform convergence on A (INT, III, p. 41, § 1, n° 1, prop. 1 and INT, III, p. 46, § 1, n° 2, prop. 5), one concludes that the linear map \(u_{k}\) is the limit of a sequence of finite-rank linear maps. It is therefore compact (cor. 1 of III, p. 4).

Consider the general case. Let \(k \in \mathcal{L}^{p', q'}(\mathrm{X}, \mathrm{Y})\) . For every \(\varepsilon > 0\) there exists a map \(k_{\varepsilon, \circ} \in \mathcal{K}(\mathrm{Y}; \mathcal{L}^{p'}(\mathrm{X}))\) such that

\[
\left\| k _ {\circ} - k _ {\varepsilon , \circ} \right\| _ {q ^ {\prime}} <   \varepsilon .
\]

The corresponding function \(k_{\varepsilon}\) : X × Y → C then satisfies

\[
\| u _ {k} - u _ {k _ {\varepsilon}} \| \leqslant \| k - k _ {\varepsilon} \| _ {p, q} \leqslant \mathrm{N} _ {p ^ {\prime}, q ^ {\prime}} (k - k _ {\varepsilon}) = \| k _ {\circ} - k _ {\varepsilon , \circ} \| _ {q ^ {\prime}} <   \varepsilon .
\]

Since \(u_{k_{\varepsilon}}\) is compact by the preceding, the map \(u_{k}\) is also compact (prop. 2 of III, p. 4).

\pandocbounded{\includegraphics[keepaspectratio]{images/part001_ec4b93bc.jpg}}

Set \(p = 1\) and suppose \(q > 1\) . There may exist mappings \(k \colon \mathrm{X} \times \mathrm{Y} \to \mathbf{C}\) such that \(k_y \in \mathcal{L}^\infty(\mathrm{X})\) for all \(y \in \mathrm{Y}\) , and such that the integral

\[
\int_ {Y} ^ {*} \| k _ {y} \| _ {\infty} ^ {q ^ {\prime}} d \nu (y)
\]

is finite (in particular, condition (2) of III, p. 27 is valid), but such that the linear mapping \(u_k\) from \(L^1(X)\) into \(L^{q'}(X)\) is not compact (exercise 2 of III, p. 119).

In particular, one deduces from the proposition the following corollary:

COROLLARY 1 (Hilbert--Schmidt). --- Let \(k \in \mathcal{L}^{2}(X \times Y)\) . The linear map \(u_{k}\) is compact from \(L^{2}(X)\) into \(L^{2}(Y)\) .

By prop. 5, one has \(k \in \mathcal{N}^{2,2}(\mathrm{X} \times \mathrm{Y})\) , so the linear map \(u_k\) is defined (prop. 3). One has \(k \in \mathcal{L}^{2,2}(\mathrm{X} \times \mathrm{Y})\) (prop. 6, \(a\) ), hence \(u_k\) is compact (prop. 7).

Remark. --- In the case p = q = 2, it is generally more convenient to express the characterization of the integral operator \(u_{k}\) using the scalar product: it is the unique map from \(\mathrm{L}^{2}(\mathrm{X})\) into \(\mathrm{L}^{2}(\mathrm{Y})\) such that

\[
\langle g \mid u _ {k} (f) \rangle = \int_ {\mathrm{X} \times \mathrm{Y}} k (x, y) f (x) \overline {{g (y)}} d (\mu \otimes \nu) (x, y)
\]

for all \(f \in \mathrm{L}^{2}(\mathrm{X})\) and \(g \in \mathrm{L}^{2}(\mathrm{Y})\) . Moreover, the adjoint of \(u_{k}\) is the linear map \(u_{k^{*}}\) , where \(k^{*} \in \mathrm{L}^{2}(\mathrm{Y} \times \mathrm{X})\) satisfies \(k^{*}(y, x) = \overline{k(x, y)}\) for almost every \((y, x) \in \mathrm{Y} \times \mathrm{X}\) . Indeed, for all \(f \in \mathrm{L}^{2}(\mathrm{X})\) and \(g \in \mathrm{L}^{2}(\mathrm{Y})\) , we have

\[
\langle g \mid u _ {k} (f) \rangle = \int_ {\mathrm{X} \times \mathrm{Y}} k (x, y) \overline {{g (y)}} f (x) d (\mu \otimes \nu) (x, y) = \langle u _ {k ^ {*}} (g) \mid f \rangle .
\]

COROLLARY 2. --- Let G be a compact topological group equipped with a Haar measure μ. Let p be a real number such that 1 \textless{} p \textless{} +∞ and let q be the conjugate exponent of p. Let \(f \in \mathcal{L}^{q}(G)\). Set k(x,y) = f(x⁻¹y) for all (x,y) ∈ G × G.

\begin{enumerate}
\def\labelenumi{\alph{enumi})}
\item
  We have \(k \in \mathcal{L}^{q,p}(\mathrm{G},\mathrm{G})\) ;
\item
  For \(\varphi \in \mathrm{L}^p (\mathrm{G})\) the convolution \(\varphi *f\) belongs to \(\mathrm{L}^p (\mathrm{G})\) and the linear map \(v_{f}\colon \varphi \mapsto \varphi *f\) of \(\mathrm{L}^p (\mathrm{G})\) into itself is continuous. It coincides with the endomorphism \(u_{k}\) . In particular, the linear map \(v_{f}\) is compact.
\end{enumerate}

Since G is compact, the function f belongs to \(\mathcal{L}^{1}(G)\) . The linear map \(v_{f}\) is therefore defined and continuous according to INT, VIII, p. 167, § 4, n° 5, prop. 13.

The map \(k\) is \((\mu \otimes \mu)\) -measurable, and we have \(\| k_y\| _q = \| f\| _q\) for all \(y\in G\) . Since the measure \(\mu\) is bounded, it follows that

\[
\int_ {\mathrm{G}} \| k _ {y} \| _ {q} ^ {p} d \mu (y) = \mu (\mathrm{G}) \| f \| _ {q} ^ {p},
\]

hence \(k \in \mathcal{L}^{q,p}(G, G)\) . The linear map \(u_{k}\) is therefore a compact linear map from \(\mathrm{L}^{p}(\mathrm{G})\) into \(\mathrm{L}^{p}(\mathrm{G})\) (prop. 7).

Let \(\varphi\in\mathcal{K}(G)\) . By INT, VIII, p. 166, § 4, no. 5, prop. 11 and prop. 4 of III, p. 29, we then have

\[
\begin{array}{c} v _ {f} (\varphi) (y) = \int_ {\mathrm{G}} \varphi (x) f (x ^ {- 1} y) d \mu (x) \\ = \int_ {\mathrm{G}} \varphi (x) k (x, y) d \mu (x) = u _ {k} (\varphi) (y) \end{array}
\]

for almost every \(y \in G\) . This implies that \(v_{f} = u_{k}\) .

